\Askhsh{20856}{2}
\begin{ylh}{1}
\item 4.3 Πολυωνυμικές εξισώσεις και ανισώσεις
\end{ylh}
Δίνεται η συνάρτηση $f(x) = 2x^{3} + x^{2} + x - 1,\ \ x \in \mathbb{R}$.
\begin{alist}
\item Να αποδείξετε οτι η εξίσωση $f(x) = 0$ δεν έχει ακέραιες ρίζες.
\monades{12}
\item Στο παρακάτω σχήμα φαίνεται η γραφική παράσταση της συνάρτησης $f$.
\end{alist}
\begin{alist}
\item
  Να δικαιολογήσετε ότι η εξίσωση $f(x) = 0$ έχει μία ρίζα.
\monades{04}
\end{alist}
\begin{alist}
\item
  Να αποδείξετε ότι η ρίζα αυτή βρίσκεται στο διάστημα $(0,1)$.
\monades{09}
% TODO: Μετατροπή σχήματος σε TikZ/PGFPlots
\begin{center}
\begin{tikzpicture}
\begin{axis}[
    axis lines = middle, xlabel = {$x$}, ylabel = {$y$},
    xmin = -3.5, xmax = 3.5, ymin = -1.5, ymax = 6.5,
    xtick = {-3,-2,-1,1,2,3}, ytick = {1,2,3,4,5,6},
    grid = both, grid style = {gray!30},
    samples = 150, width = 10cm, height = 7cm,
    tick label style = {font=\small},
]
    \addplot[ultra thick, maincolor, domain=-3:1.8] {exp(x)} node[above left, pos=0.9] {$C_f$};
    \filldraw (axis cs:0,1) circle (2pt) node[above right] {$1$};
    \filldraw (axis cs:1,{exp(1)}) circle (2pt) node[right] {$A$};
    \draw[dashed] (axis cs:1,{exp(1)}) -- (axis cs:1,0) node[below] {$1$};
\end{axis}
\end{tikzpicture}
% % \caption{Σχήμα 1}
\end{center}
\end{alist}

